% --- CAPÍTULO III: MARCO METODOLÓGICO ---
\chapter{Marco Metodológico}
\chaptermark{Capítulo III: Marco Metodológico}

\section{Nivel y tipo de investigación}
\lipsum[18]
De acuerdo con los lineamientos metodológicos para trabajos de grado en Venezuela, es crucial establecer con claridad el nivel y tipo de investigación que sustentan el estudio \autocite{book:hernandez2014}.

\section{Diseño de la investigación}
\lipsum[19]
El diseño de investigación adoptado es no experimental, transversal de alcance descriptivo, siguiendo los preceptos establecidos para este tipo de estudios en el área de ingeniería \parencite{book:hernandez2014}.

\section{Población y muestra}
\lipsum[20]
La población estuvo conformada por 450 individuos. La selección de la muestra se realizó mediante un muestreo probabilístico, tal como se sugiere en los textos fundamentales de metodología.

\section{Técnicas e instrumentos de recolección de datos}
\lipsum[21]
Para la recolección de los datos, se empleó la técnica de la encuesta y se utilizó el cuestionario como instrumento. Este instrumento fue validado por tres expertos en el área, siguiendo el proceso descrito por la \textcite{online:unefa2022}. El cuestionario fue realizado con \Gls{latex} y \gls{maths}.

\section{Técnicas de procesamiento y análisis de datos}
\lipsum[22]
Los datos cuantitativos obtenidos se procesaron utilizando el *Software Estadístico para Ciencias Sociales* (SPSS). El análisis se centró en estadísticas descriptivas y la aplicación de modelos estadísticos no paramétricos para la comparación de grupos \parencite{article:perez2020}.

\section{Cronograma de Actividades}
El cronograma de actividades se presenta a continuación, detallando la duración estimada de cada fase del proyecto en semanas, y se visualiza en la Tabla \ref{tab:gantt}. 
% ¡IMPORTANTE! Asegúrese de que el texto de su sección ahora dice: 
% "...se visualiza en la Tabla \ref{tab:gantt}." 
% para que siempre se muestre el número correcto.

\clearpage % Salta a la siguiente página antes de rotar

\begin{landscape} % <-- INICIO DE LA PÁGINA HORIZONTAL
\begin{table}[t] % Entorno de tabla que genera el número 3.X e indexa
\centering
% ------------------------------------------------
% CAPTION & LABEL (SOLUCIÓN AL NÚMERO Y AL ÍNDICE)
% ------------------------------------------------
\caption{Cronograma de actividades del proyecto}
\label{tab:gantt}
\vspace{0.3cm} % Espacio entre el título y el gráfico

\begin{ganttchart}[
    % CONFIGURACIÓN GLOBAL
    hgrid,
    vgrid={*4{dotted}, *4{dotted}, *5{dotted}},
    x unit=1.0cm, % Aprovechamos el ancho extra de la hoja horizontal
    y unit chart=0.7cm,
    % ESTILOS VISUALES
    title/.style={fill=gray!30, font=\bfseries\normalsize},
    title label font=\bfseries, 
    bar/.style={fill=blue!40, draw=blue!60, rounded corners=2pt},
    bar label font=\normalsize,
    milestone/.style={fill=red!60, draw=red!70, rounded corners=2pt},
    group/.style={draw=black, fill=gray!15},
    link/.style={->, thick} 
    ]{1}{13}
    
    % ESTRUCTURA DE TIEMPO
    \gantttitle{Periodo de Ejecución (Semanas)}{13} \\
    \gantttitle{Noviembre}{4}
    \gantttitle{Diciembre}{4}
    \gantttitle{Enero}{5} \\
    \gantttitlelist{1,...,13}{1} \\
    
    % BARRAS DE ACTIVIDADES (Con nombres explícitos para el link)
    \ganttbar[name=elem0]{Determinar requerimientos del sistema}{1}{4} \\
    \ganttbar[name=elem1]{Elaborar diseño del sistema electrónico}{3}{7} \\
    \ganttbar[name=elem2]{Realizar programas de aplicación}{5}{12} \\
    \ganttbar[name=elem3]{Validación del sistema}{10}{13} \\
    
    % HITOS
    \ganttmilestone{Requerimientos aprobados}{4} \\
    \ganttmilestone{Diseño completado}{7} \\
    \ganttmilestone{Programación terminada}{12} \\
    \ganttmilestone{Validación exitosa}{13}
    
    % DEPENDENCIAS
    \ganttlink{elem0}{elem1}
    \ganttlink{elem1}{elem2}
    \ganttlink{elem2}{elem3}
    
\end{ganttchart}

\end{table}
\end{landscape} % <-- FIN DE LA PÁGINA HORIZONTAL

\clearpage % Regresa al modo vertical para el resto del documento